\documentclass[../readings.tex]{subfiles}

\begin{document}

\subsection{Conjugate Gradient Method}
\label{sub:cg}

\textBox{%
The Conjugate Gradient (CG) method is a standard iterative solver for large, sparse, symmetric positive definite (SPD) linear systems. It belongs to the class of \textbf{Krylov subspace methods}, which seek an approximate solution within a space spanned by powers of the matrix applied to the initial residual.
}

\addRemark{%
\textbf{(Assumption)}\enspace CG requires $\bA$ to be symmetric positive definite. If $\bA$ is nonsymmetric, use GMRES or another nonsymmetric Krylov method. If $\bA$ is symmetric indefinite, use methods designed for indefinite systems, such as MINRES.
}

\textBox{%
For SPD $\bA$, solving $\bA\bx=\bb$ is equivalent to minimizing a convex quadratic energy. CG chooses search directions that do not interfere with previous progress, so each step is optimal over a growing Krylov subspace.
}

\subsubsection{Krylov Subspaces}

\addDef{Krylov Subspace}{%
The $k$-th Krylov subspace $\mathcal{K}_k(\bA, \bb)$ is defined as the span of the first $k$ vectors in the sequence $\{\bb, \bA\bb, \bA^2\bb, \dots \}$:
\begin{align*}
    \mathcal{K}_k(\bA, \bb) = \text{span}\{ \bb, \bA\bb, \bA^2\bb, \dots, \bA^{k-1}\bb \}.
\end{align*}
Krylov methods are efficient because they only require the ability to compute matrix-vector products $\bA\bv$, making them ideal for sparse matrices where $\bA$ is never explicitly stored.
}
\label{def:krylov-subspace-cg}

\subsubsection{CG as a Projection Method}

\textBox{%
For SPD $\bA$, the quadratic
\begin{align*}
    \phi(\bx)=\frac{1}{2}\bx^T\bA\bx-\bb^T\bx
\end{align*}
has gradient $\nabla\phi(\bx)=\bA\bx-\bb$. Therefore $\nabla\phi(\bx)=0$ exactly when $\bA\bx=\bb$. Solving the linear system is the same as finding the unique minimizer of $\phi$.
}

\addRemark{%
\textbf{Error Minimization Properties:}\enspace
For an SPD system $\bA\bx = \bb$, the CG method identifies the unique vector $\bx_k \in \mathcal{K}_k(\bA, \bb)$ that minimizes the \textbf{$\bA$-norm of the error}:
\begin{align*}
    \bx_k = \arg\min_{\by \in \mathcal{K}_k} \|\bx^* - \by\|_\bA,
\end{align*}
where $\bx^*$ is the exact solution and $\|\be\|_\bA = \sqrt{\be^T\bA\be}$. This is equivalent to minimizing the quadratic energy functional $\phi(\by) = \frac{1}{2} \by^T \bA \by - \by^T \bb$.
}

\subsubsection{Conjugate Directions and A-Orthogonality}

\addRemark{%
\textbf{The Basis of Directions:}\enspace
To find this minimum efficiently, CG constructs a basis of search directions $\{\bp_0, \bp_1, \dots, \bp_{k-1}\}$ that are \textbf{$\bA$-orthogonal} (or conjugate), meaning $\bp_i^T \bA \bp_j = 0$ for $i \neq j$.
\begin{itemize}
    \item \textbf{Global Optimality:}\enspace Because directions are $\bA$-orthogonal, the $k$-th update $\bx_k = \bx_{k-1} + \alpha_{k-1}\bp_{k-1}$ is guaranteed to be the best possible update in the entire subspace $\mathcal{K}_k$. There is no need to revisit previous directions.
    \item \textbf{Short Recurrence:}\enspace Due to the symmetry of $\bA$, the new conjugate direction $\bp_k$ can be computed using only the current residual and the \emph{single} previous direction. This leads to the $O(n)$ storage requirement.
\end{itemize}
}

\addRemark{%
\textbf{(Residuals vs.\ directions)}\enspace In exact arithmetic, CG residuals are mutually orthogonal in the Euclidean inner product, while search directions are mutually orthogonal in the $\bA$-inner product. These two orthogonality structures are what make the short recurrence possible.
}

\addRemark{%
\textbf{Why not steepest descent?}\enspace Steepest descent always moves in the negative gradient direction, but successive gradients can undo previous progress by zig-zagging across narrow energy valleys. CG corrects this by choosing conjugate directions adapted to the geometry of $\bA$.
}

\subsubsection{The CG Algorithm}

\addDef{Conjugate Gradient Algorithm}{%
\textbf{Initialize:}\enspace $\bx^{(0)} = \bzero, \br^{(0)} = \bb, \bp^{(0)} = \br^{(0)}$.\\
\textbf{Iteration:}\enspace
\begin{enumerate}
    \item \textbf{$\alpha_k = \frac{\br^{(k)T}\br^{(k)}}{\bp^{(k)T}\bA\bp^{(k)}}$} (Step length)
    \item \textbf{$\bx^{(k+1)} = \bx^{(k)} + \alpha_k \bp^{(k)}$} (Update iterate)
    \item \textbf{$\br^{(k+1)} = \br^{(k)} - \alpha_k \bA\bp^{(k)}$} (Update residual)
    \item \textbf{$\beta_k = \frac{\br^{(k+1)T}\br^{(k+1)}}{\br^{(k)T}\br^{(k)}}$} (Gram-Schmidt weight)
    \item \textbf{$\bp^{(k+1)} = \br^{(k+1)} + \beta_k \bp^{(k)}$} (Next direction)
\end{enumerate}
}
\label{alg:cg}

\addRemark{%
\textbf{Efficiency:}\enspace
\begin{itemize}
    \item \textbf{Matvecs:}\enspace Only 1 matrix-vector product per iteration.
    \item \textbf{Memory:}\enspace $O(n)$ storage (requires only a few vectors). No need to store the basis.
\end{itemize}
}

\addRemark{%
\textbf{(Practical stopping)}\enspace A common stopping rule is
\begin{align*}
    \frac{\|\br^{(k)}\|_2}{\|\bb\|_2} \leq \texttt{tol}.
\end{align*}
For ill-conditioned systems, residual decrease may not translate directly into error decrease. The condition number controls that conversion.
}

\subsubsection{Convergence Rate}

\addThm{Convergence Bound}{%
\begin{align*}
    \frac{\|\be^{(k)}\|_\bA}{\|\be^{(0)}\|_\bA} \leq 2 \left( \frac{\sqrt{\kappa}-1}{\sqrt{\kappa}+1} \right)^k.
\end{align*}
}
\label{thm:cg-convergence}

\addRemark{%
\textbf{The $\sqrt{\kappa}$ Advantage:}\enspace
While stationary methods depend on $(\kappa-1)/(\kappa+1)$, CG depends on the \textbf{square root} of the condition number. For $\kappa=100$, CG reduces error $5\times$ faster per step.
}

\addRemark{%
\textbf{Superlinear Convergence:}\enspace
If eigenvalues are clustered, CG converges much faster than the worst-case bound suggests (e.g., $k \ll n$ for clustered spectra).
}

\addRemark{%
\textbf{(Polynomial view)}\enspace CG chooses a degree-$k$ polynomial $p_k$ with $p_k(0)=1$ so that
\begin{align*}
    \be^{(k)} = p_k(\bA)\be^{(0)}
\end{align*}
is small in the $\bA$-norm. Good convergence occurs when a low-degree polynomial can be small on the eigenvalues of $\bA$.
}
\subsubsection{Preconditioning}

The convergence of CG depends on the condition number $\kappa(\bA)$. \textbf{Preconditioning} transforms the system using a symmetric positive definite matrix $\bM$ that approximates $\bA$ but is much easier to invert.

\addRemark{%
\textbf{Spectral Transformation:}\enspace
The goal of the preconditioner is to cluster the eigenvalues of $\bM^{-1}\bA$ near 1. A perfectly preconditioned system ($\bM = \bA$) has all eigenvalues equal to 1 and converges in a single step. In practice, a good preconditioner reduces the effective condition number, significantly accelerating convergence.
}

\addRemark{%
\textbf{(Preconditioning is not optional in hard problems)}\enspace For large PDE systems, unpreconditioned CG can require too many iterations. The preconditioner is often the main algorithmic design choice.
}

\addDef{Common Preconditioners}{%
\begin{enumerate}
    \item \textbf{Jacobi:}\enspace $\bM = \text{diag}(\bA)$. (Cheap, moderate speedup).
    \item \textbf{Incomplete Cholesky (IC):}\enspace $\bM \approx \bL\bL^T$ with sparse factors. (Robust).
    \item \textbf{Multigrid:}\enspace Near $O(n)$ total complexity for PDE systems.
\end{enumerate}
}
\label{def:preconditioners}

\addExcr{%
\begin{enumerate}
    \item \textbf{Convergence Check:}\enspace Solve a $100 \times 100$ Poisson system via CG. Compare your error plot to the theoretical bound.
    \item \textbf{Spectral Clustering:}\enspace Construct $\bA$ with eigenvalues $\{1, \dots, 1, 100, 200\}$. Observe superlinear convergence in the first few steps.
    \item \textbf{Preconditioning:}\enspace Compare iteration counts for a sparse system with and without an IC preconditioner.
\end{enumerate}
}

\end{document}
